\chaptertitle{Aria with Diverse Variations}

                     Aria with Diverse Variations
     Achilles has been unable to 'sleep these past few nights. His friend the
     Tortoise has come over tonight, to keep him company during these annoying
     hours.

Tortoise: I am so sorry to hear of the troubles that have been plaguing you, my dear
    Achilles. I hope my company will provide a welcome relief from all the unbearable
    stimulation which has kept you awake. Perhaps I will bore you sufficiently that you
    will at long last go to sleep. In that way, I will be of some service.
Achilles: Oh, no, I am afraid that I have already had some of the world's finest bores try
    their hand at boring me to sleep-and all, sad to say, to no avail. So you will be no
    match for them. No, Mr. T, I invited you over hoping that perhaps you could
    entertain me with a little this or that, taken from number theory, so that I could at
    least while away these long hours in an agreeable fashion. You see, I have found that
    a little number theory does wonders for my troubled psyche.
Tortoise: How quaint an idea! You know, it reminds me, just a wee bit, of the story of
    poor Count Kaiserling.
'Achilles: Who was he?
Tortoise: Oh, he was a Count in Saxony in the eighteenth century-a Count of no account,
    to tell the truth-but because of him-well, shall I tell you the story? It is quite
    entertaining.
Achilles: In that case, by all means, do!
Tortoise: There was a time when the good Count was suffering from sleeplessness, and it
    just so happened that a competent musician lived in the same town, and so Count
    Kaiserling commissioned this musician to compose a set of variations to be played
    by the Count's court harpsichordist for him during his sleepless nights, to make the
    hours pass by more pleasantly.
Achilles: Was the local composer up to the challenge?
Tortoise: I suppose so, for after they were done, the Count rewarded him most
    lucratively-he presented him with a gold goblet containing one hundred Louis d'or.
Achilles: You don't say! I wonder where he came upon such a goblet and all those Louis
    d'or, in the first place.
Tortoise. Perhaps he saw it in a museum, and took a fancy to it.
Achilles: Are you suggesting he absconded with it?
Tortoise: Now, now, I wouldn't put it exactly that way, but ... Those days, Counts could
    get away with most anything. Anyway, it is clear that the Count was most pleased
    with the music, for he was constantly entreating his harpsichordist-a mere lad of a
    fellow, name of Goldberg-to




Aria with Diverse Variations                                                          397

    play one or another of these thirty variations. Consequently (and somewhat
    ironically) the variations became attached to the name of young Goldberg, rather
    than to the distinguished Count's name.
Achilles: You mean, the composer was Bach, and these were the so-called "Goldberg
    Variations"?
Tortoise: Do I ever! Actually, the work was entitled Aria with Diverse Variations, of
    which there are thirty. Do you know how Bach structured these thirty magnificent
    variations?
Achilles: Do tell.
Tortoise: All the pieces-except the final one-are based on a single theme, which he called
    an "aria". Actually, what binds them all together is not a common melody, but a
    common harmonic ground. The melodies may vary, but underneath, there is a
    constant theme. Only in the last variation did Bach take liberties. It is a sort of "post-
    ending ending". It contains extraneous musical ideas having little to do with the
    original Theme-in fact, two German folk tunes. That variation is called a
    "quodlibet".
Achilles: What else is unusual about the Goldberg Variations?
Tortoise: Well, every third variation is a canon. First a canon in two canonizing voices
    enter on the SAME note. Second, a canon in which one of the canonizing voices
    enters ONE NOTE HIGHER than the first. Third, one voice enters Two notes higher
    than the other. And so on, until the final canon has entries just exactly one ninth
    apart. Ten canons, all told. And
Achilles: Wait a minute. Don't I recall reading somewhere or other about fourteen
    recently discovered Goldberg canons ...
Tortoise: Didn't that appear in the same journal where they recently reported the
    discovery of fourteen previously unknown days in November?
Achilles: No, it's true. A fellow named Wolff-a musicologist-heard about a special copy
    of the Goldberg Variations in Strasbourg. He went there to examine it, and to his
    surprise, on the back page, as a sort of "post-ending ending", he found these fourteen
    new canons, all based on the first eight notes of the theme of the Goldberg
    Variations. So now it is known that there are in reality forty-four Goldberg
    Variations, not thirty.
Tortoise: That is, there are forty-four of them, unless some other musicologist discovers
    yet another batch of them in some unlikely spot. And although it seems improbable,
    it is still possible, even if unlikely, that still another batch will be discovered, and
    then another one, and on and on and on ... Why, it might never stop! We may never
    know if or when we have the full complement of Goldberg Variations.
Achilles: That is a peculiar idea. Presumably, everybody thinks that this latest discovery
    was just a fluke, and that we now really do have all the Goldberg Variations. But just
    supposing that you are right, and some turn up sometime, we shall start to expect
    this kind of thing. At




Aria with Diverse Variations                                                              398

    that point, the name "Goldberg Variations" will start to shift slightly in meaning, to
    include not only the known ones, but also any others which might eventually turn up.
    Their number-call it 'g'-is certain to be finite, wouldn't you agree?-but merely
    knowing that g is finite isn't the same as knowing how big g is. Consequently, this
    information won't tell us when the last Goldberg Variation has been located.
Tortoise: That is certainly true.
Achilles: Tell me-when was it that Bach wrote these celebrated variations?
Tortoise: It all happened in the year 1742, when he was Cantor in Leipzig. Achilles:
    1742? Hmm ... That number rings a bell.
Tortoise: It ought to, for it happens to be a rather interesting number, being a sum of two
    odd primes: 1729 and 13.
Achilles: By thunder! What a curious fact' I wonder how often one runs across an even
    number with that property. Let's see

                       6= 3+3
                       8= 3+5
                       10= 3+7=      5+5
                       12= 5+7
                       14=3+11= 7+7
                       16=3+ 13= 5+ 11
                       18=5+ 13= 7+ 11
                       20=3+ 17= 7+ 13
                       22 = 3 + 19 = 5 + 17 =      11       + 11
                       24=5+19= 7+17=11+13
                       26=3+23= 7+19=13+13
                       28 = 5 + 23 = 11 + 17
                       30 = 7 + 23 = 11 + 19 = 13 + 17

    Now what do you know-according to my little table here, it seems to be quite a
    common occurrence. Yet I don't discern any simple regularity in the table so far.

Tortoise: Perhaps there is no regularity to be discerned.
Achilles: But of course there is! I am just not clever enough to spot it right off the bat.
Tortoise: You seem quite convinced of it.
Achilles: There's no doubt in my mind. I wonder ... Could it be that ALL even numbers
    (except 4) can be written as a sum of two odd primes?
Tortoise: Hmm ... That question rings a bell ... Ah, I know why! You're not the first
    person to ask that question. Why, as a matter of fact, in the year 1742, a
    mathematical amateur put forth this very question in a
Achilles: Did you say 1742? Excuse me for interrupting, but I just noticed that 1742
    happens to be a rather interesting number, being a difference of two odd primes:
    1747 and 5.
Tortoise: By thunder! What a curious fact! I wonder how often one runs across an even
    number with that property.




Aria with Diverse Variations                                                           399

Achilles: But please don't let me distract you from your story.
Tortoise: Oh, yes-as I was saying, in 1742. a certain mathematical amateur, whose name
    escapes me momentarily, sent a letter to Euler, who at the time was at the court of
    King Frederick the Great in Potsdam, and-well, shall I tell you the story? It is not
    without charm. Achilles: In that case, by all means, do!
Tortoise: Very well. In his letter, this dabbler in number theory propounded an unproved
    conjecture to the great Euler: "Every even number can he represented as a sum of
    two odd primes." Now what was that fellow's name?
Achilles: I vaguely recollect the story, from some number theory book or other. Wasn't
    the fellow named Iiupfergiidel
Tortoise: Hmm ... No, that sounds too long.
Achilles: Could it have been "Silberescher"?
Tortoise: No, that's not it, either. There's a name on the tip of' my tongue-ah-ah-oh yes! It
    was "Goldbach"! Goldbach was the fellow. Achilles: I knew it was something like
    that.
Tortoise: Yes-your guesses helped jog my memory. It's quite odd, how one occasionally
    has to hunt around in one's memory as if for a book in a library without call numbers
    ... But let us get back to 1742.
Achilles: Indeed, let's. I wanted to ask you: did Euler ever prove that this guess by
    Goldbach was right?
Tortoise: Curiously enough, he never even considered it worthwhile working on.
    However, his disdain was not shared by all mathematicians. In fact, it caught the
    fancy of many, and became known as the "Goldbach Conjecture".
Achilles: Has it ever been proven correct?
Tortoise: No, it hasn't. But there have been some remarkable near misses. For instance, in
    1931 the Russian number theorist Schnirelmann proved that any number-even or
    odd-can be represented as the sum of not more than 300,000 primes.
Achilles: What a strange result. Of what good is it?
Tortoise: It has brought the problem into the domain of the finite. Previous to
    Schnirelmann's proof, it was conceivable that as you took larger and larger even
    numbers, they would require more and more primes to represent them. Some even
    number might take a trillion primes to represent it! Now it is known that that is not
    so-a sum of 300,000 primes (or fewer) will always suffice.
Achilles: I see.
Tortoise: Then in 1937, a sly fellow named Vinogradov-a Russian too-managed to
    establish something far closer to the desired result: namely, every sufficiently large
    ODD number can be represented as a sum of no more than THREE odd primes. For
    example, 1937 = 641 + 643 + 653. We could say that an odd number which is
    representable as a sum of three odd primes has "the Vinogradov property. Thus, all
    sufficiently large odd numbers have the Vinigradov properties




Aria with Diverse Variations                                                             400

Achilles: Very well-but what does "sufficiently large" mean?
Tortoise: It means that some finite number of odd numbers may fail to have the
    Vinogradov property, but there is a number-call it 'v'beyond which all odd numbers
    have the Vinogradov property. But Vinogradov was unable to say how big v is. So in
    a way, v is like g, the finite but unknown number of Goldberg Variations. Merely
    knowing that v is finite isn't the same as knowing how big v is. Consequently, this
    information won't tell us when the last odd number which needs more than three
    primes to represent it has been located.
Achilles: I see. And so any sufficiently large even number 2N can be represented as a
    sum of FOUR primes, by first representing 2N - 3 as a sum of three primes, and then
    adding back the prime number 3.
Tortoise: Precisely. Another close approach is contained in the Theorem which says, "All
    even numbers can be represented as a sum of one prime and one number which is a
    product of at most two primes."
Achilles: This question about sums of two primes certainly leads you into strange
    territory. I wonder where you would be led if you looked at DIFFERENCES of two
    odd primes. I'll bet I could glean some insight into this teaser by making a little table
    of even numbers, and their representations as differences of two odd primes, just as I
    did for sums. Let's see ...

         2= 5-3,        7-5,          13-11,          19-17, etc.
         4 = 7 - 3,     11 - 7,       17 - 13,        23 - 19,etc.
         6 = 11 - 5,    13 - 7,       17 - 11,        19- 13, etc.
         8 = 11 - 3,    13 - 5,       19 - 11,        31 - 23,etc.
         10 = 13 - 3,   17 - 7,       23 - 13,        29- 19, etc.

My gracious! There seems to be no end to the number of different representations I can
    find for these even numbers. Yet I don't discern any simple regularity in the table so
    far.
Tortoise: Perhaps there is no regularity to be discerned.
Achilles: Oh, you and your constant rumblings about chaos! I'll hear none of that, thank
    you.
Tortoise: Do you suppose that EVERY even number can be represented somehow as the
    difference of two odd primes?
Achilles: The answer certainly would appear to be yes, from my table. But then again, I
    suppose it could also be no. That doesn't really get us very far, does it?
Tortoise: With all due respect, I would say there are deeper insights to be had on the
    matter.
Achilles: Curious how similar this problem is to Goldbach's original one. Perhaps it
    should be called a "Goldbach Variation".
Tortoise: Indeed. But you know, there is a rather striking difference between the
    Goldbach Conjecture, and this Goldbach Variation, which I would like to tell you
    about. Let us say that an even number 2N has the “Goldbach property” if it is the
    SUM of two odd primes, and it has the “Tortoise property” if it is the DIFFERENCE
    of two odd primes

Aria with Diverse Variations                                                             401

Achilles: I think you should call it the "Achilles property". After all, I suggested the
    problem.
Tortoise: I was just about to propose that we should say that a number which LACKS the
    Tortoise property has the "Achilles property". Achilles: Well, all right .. .
Tortoise: Now consider, for instance, whether I trillion has the Goldbach property or the
    Tortoise property. Of course, it may have both.
Achilles: I can consider it, but I doubt whether I can give you an answer to either
    question.
Tortoise: Don't give up so soon. Suppose I asked you to answer one or the other question.
    Which one would you pick to work on?
Achilles: I suppose I would flip a coin. I don't see much difference between them.
Tortoise: Aha: But there's a world of difference' If you pick the Goldbach property,
    involving SUMS of primes, then you are limited to using primes which are bounded
    between 2 and 1 trillion, right?
Achilles: Of course.
Tortoise: So your search for a representation for 1 trillion as a sum of two primes is
    GUARANTEED TO TERMINATE.
Achilles: Ahhh! I see your point. Whereas if I chose to work on representing 1 trillion as
    the DIFFERENCE of two primes, I would not have any bound on the size of the
    primes involved. They might be so big that it would take me a trillion years to find
    them.
Tortoise: Or then again, they might not even EXIST! After all, that's what the question
    was asking-do such primes exist, It wasn't of much concern how big they might turn
    out to be.
Achilles: You're right. If they didn't exist, then a search process would lead on forever,
    never answering yes, and never answering no. And nevertheless, the answer would
    be no.
Tortoise: So if you have some number, and you wish to test whether it has the Goldbach
    property or the Tortoise property, the difference between the two tests will be this: in
    the former, the search involved is GUARANTEED TO TERMINATE; in the latter,
    it is POTENTIALLY ENDLESS-there are no guarantees of any type. It might just
    go merrily on forever, without yielding an answer. And yet, on the other hand, in
    some cases, it might stop on the first step.
Achilles: I see there is a rather vast difference between the Goldbach and Tortoise
    properties.
Tortoise: Yes, the two similar problems concern these vastly different properties. The
    Goldbach Conjecture is to the effect that all even numbers have the Goldbach
    property; the Goldbach Variation suggests that all even numbers have the Tortoise
    property. Both problems are unsolved, but what is interesting is that although they
    sound very much alike, they involve properties of whole numbers which are quite
    different.
Achilles: I see what you mean. The Goldbach property is a detectable, or




Aria with Diverse Variations                                                            402

    recognizable property of any even number, since I know how to test for its presence
    just embark on a search. It will automatically come to an end with a yes or no
    answer. The Tortoise property, however, is more elusive, since a brute force search
    just may never give an answer.
Tortoise: Well, there may be cleverer ways of searching in the case of the Tortoise
    property, and maybe following one of them would always come to an end, and yield
    an answer.
Achilles: Couldn't the search only end if the answer were yes\%
Tortoise: Not necessarily. There might be some way of proving that whenever the search
    lasts longer than a certain length of time, then the answer must be no. There might
    even be some OTHER way of searching for the primes, not such a brute force way,
    which is guaranteed to find them if they exist, and to tell if they don't. In either case,
    a finite search would be able to yield the answer no. But I don't know if such a thing
    can be proven or not. Searching through infinite spaces is always a tricky matter, you
    know.
Achilles: So as things stand now, you know of no test for the Tortoise property which is
    guaranteed to terminate-and yet there MIGHT exist such a search.
Tortoise: Right. I suppose one could embark on a search for such a search, but I can give
    no guarantee that that "meta-search" would terminate, either.
Achilles: You know, it strikes me as quite peculiar that if some even number-for
    example, a trillion-failed to have the Tortoise property, it would be caused by an
    infinite number of separate pieces of information. It's funny to think of wrapping all
    that information up into one bundle, and calling it, as you so gallantly suggested,
    "the Achilles property" of 1 trillion. It is really a property of the number system as a
    "HOLE, not just of the number 1 trillion.
Tortoise: That is an interesting observation, Achilles, but I maintain that it makes a good
    deal of sense to attach this fact to the number 1 trillion nevertheless. For purposes of
    illustration, let me suggest that you . consider the simpler statement "29 is prime".
    Now in fact, this statement really means that 2 times 2 is not 29, and 5 times 6 is not
    29, and so forth, doesn't it?
Achilles: It must, I suppose.
Tortoise: But you are perfectly happy to collect all such facts together, and attach them in
    a bundle to the number 29, saying merely, "29 is prime"
Achilles: Yes ...
Tortoise: And the number of facts involved is actually infinite, isn't it:, After all, such
    facts as "4444 times 3333 is not 29" are all part of it, aren't they\%
Achilles: Strictly speaking, I suppose so. But you and I both know that you can't produce
    29 by multiplying two numbers which are both bigger than 29. So in reality, saying
    "29 is prime" is only summarizing a FINITE number of facts about multiplication




Aria with Diverse Variations                                                              403

Tortoise: You can put it that way if you want, but think of this: the fact that two numbers
    which are bigger than 29 can't have a product equal to 29 involves the entire
    structure of the number system. In that sense, that fact in itself is a summary of an
    infinite number of facts. You can't get away from the fact, Achilles, that when you
    say "29 is prime'-', you are actually stating an infinite number of things.
Achilles: Maybe so, but it feels like just one fact to me.
Tortoise: That's because an infinitude of facts are contained in your prior knowledge-they
    are embedded implicitly in the way you visualize things. You don't see an explicit
    infinity because it is captured implicitly inside the images you manipulate.
Achilles: I guess that you're right. It still seems odd to lump a property of the entire
    number system into a unit, and label the unit "primeness of 29"
Tortoise: Perhaps it seems odd, but it is also quite a convenient way to look at things.
    Now let us come back to your hypothetical idea. If, as you suggested, the number 1
    trillion has the Achilles property, then no matter what prime you add to it, you do not
    get another prime. Such a state of affairs would be caused by an infinite number of
    separate mathematical "events". Now do all these "events" necessarily spring from
    the same source? Do they have to have a common cause? Because if they don't, then
    some sort of "infinite coincidence" has created the fact, rather than an underlying
    regularity.
Achilles: An "infinite coincidence"? Among the natural numbers, NoTHING is
    coincidental-nothing happens without there being some underlying pattern. Take 7,
    instead of a trillion. I can deal with it more easily, because it is smaller. 7 has the
    Achilles property.
Tortoise: You're sure?
Achilles: Yes. Here's why. If you add 2 to it, you get 9, which isn't prime. And if you add
    any other prime to 7, you are adding two odd numbers, resulting in an even number-
    thus you again fail to get a prime. So here the "Achilleanity" of 7, to coin a term, is a
    consequence of just Two reasons: a far cry from any "infinite coincidence". Which
    just goes to support my assertion: that it never takes an infinite number of reasons to
    account for some arithmetical truth. If there WERE some arithmetical fact which
    were caused by an infinite collection of unrelated coincidences, then you could never
    give a finite proof for that truth. And that is ridiculous.
Tortoise: That is a reasonable opinion, and you are in good company in making it.
    However
Achilles: Are there actually those who disagree with this view? Such people would have
    to believe that there are "infinite coincidences", that there is chaos in the midst of the
    most perfect, harmonious, and beautiful of all creations: the system of natural
    numbers.
Tortoise: Perhaps they do; but have you ever considered that such chaos might be an
    integral part of the beauty and harmony?




Aria with Diverse Variations                                                              404

FIGURE 71. Order and Chaos, by M. C. Escher (lithograph, 1950).

Achilles: Chaos, part of perfection? Order and chaos make a pleasing unity? Heresy!
Tortoise: Your favorite artist, M. C. Escher, has been known to suggest such a heretical
    point of view in one of his pictures ... And while we're on the subject of chaos, I
    believe that you might be interested in hearing about two different categories of
    search, both of which are guaranteed to terminate.
Achilles: Certainly.




Aria with Diverse Variations                                                        405

Tortoise: The first type of search-the non-chaotic type-is exemplified by the test involved
    in checking for the Goldbach property. You just look at primes less than 2N, and if
    some pair adds up to 2N, then 2N has the Goldbach property; otherwise, it doesn't.
    This kind of test is not only sure to terminate, but you can predict BY "'HEN it will
    terminate, as well.
Achilles: So it is a PREDICTABLY TERMINATING test. Are you going to tell me that
    checking for some number-theoretical properties involves tests which are guaranteed
    to terminate, but about which there is no way to know in advance how long they will
    take?
Tortoise: How prophetic of you, Achilles. And the existence of such tests shows that
    there is intrinsic chaos, in a certain sense, in the natural number system.
Achilles: Well, in that case, I would have to say that people just don't know enough about
    the test. If they did a little more research, they could figure out how long it will take,
    at most, before it terminates. After all, there must always be some rhyme or reason to
    the patterns among integers. There can't just be chaotic patterns which defy
    prediction'
Tortoise: I can understand your intuitive faith, Achilles. However, it's not always
    justified. Of course, in many cases you are exactly right just because somebody
    doesn't know something, one can't conclude that it is unknowable' But there are
    certain properties of integers for which terminating tests can be proven to exist, and
    yet about which it can also be PROVEN that there is no way to predict in advance
    how long they will take.
Achilles: I can hardly believe that. It sounds as if the devil himself managed to sneak in
    and throw a monkey wrench into God's beautiful realm of natural numbers'
Tortoise: Perhaps it will comfort you to know that it is by no means easy, or natural, to
    define a property for which there is a terminating but not PREDICTABLY
    terminating test. Most "natural" properties of integers do admit of predictably
    terminating tests. For example, primeness. squareness, being a power of ten, and so
    on.
Achilles: Yes, I can see that those properties are completely straightforward to test for.
    Will you tell me a property for which the only possible test is a terminating but
    nonpredictable one?
Tortoise: That's too complicated for me in my sleepy state. Let me instead show you a
    property which is very easy to define, and yet for which no terminating test is
    known. I'm not saying there won't ever be one discovered, mind you just that none is
    known. You begin with a number-would you care to pick one?
Achilles: How about 15?
Tortoise: An excellent choice. We begin with your number, and if it is ODD, we triple it,
    and add 1. If it is EVEN, we take half of it. Then we repeat the process. Call a
    number which eventually reaches 1 this way a WONDROUS number, and a number
    which doesn't, an UNWONDROUS number




Aria with Diverse Variations                                                              406

Achilles: Is 15 wondrous, or unwondrous? Let's see:

15 is ODD, so I make 3n + 1:                  46
46 is EVEN, so I take half:            23
23 is ODD, so I make 3n + 1:           70
70 is EVEN, so I take half:            35
35 is ODD, so I make 3n + 1:           106
106 is EVEN, so I take half:           53
53 is ODD, so I make 3n + 1:           160
160 is EVEN, so I take half:           80
80 is EVEN, so I take half:            40
40 is EVEN, so I take half:            20
20 is EVEN, so I take half:            10
10 is EVEN, so I take half:            5
5 is ODD, so I make 3n + 1:            16
16 is EVEN, so I take half:            8
8 is EVEN, so I take half:             4
4 is EVEN, so I take half:             2
2 is EVEN, so I take half:             1
.
    Wow! That's quite a roundabout journey, from 15 to 1. But I finally reached it. That
    shows that 15 has the property of being wondrous. I wonder what numbers are
    UNwondrous ...
Tortoise: Did you notice how the numbers swung up and down, in this simply defined
    process?
Achilles: Yes. I was particularly surprised, after thirteen turns, to find myself at 16, only
    one greater than 15, the number I started with. In one sense, I was almost back where
    I started-yet in another sense, I' was nowhere near where I had started. Also, I found
    it quite curious that I had to go as high as 160 to resolve the question. I wonder how
    come.
Tortoise: Yes, there is an infinite "sky" into which you can sail, and it is very hard to
    know in advance how high into the sky you will wind up sailing. Indeed, it is quite
    plausible that you might just sail up and up and up, and never come down.
Achilles: Really? I guess that is conceivable-but what a weird coincidence it would
    require! You'd just have to hit odd number after odd number, with only a few evens
    mixed in. I doubt if that would ever happen-but I just don't know for sure.
Tortoise: Why don't you try starting with 27? Mind you, I don't promise anything. But
    sometime, just try it, for your amusement. And I'd advise you to bring along a rather
    large sheet of paper.
Achilles: Hmm ... Sounds interesting. You know, it still makes me feel funny to associate
    the wondrousness (or unwondrousness) with the starting number, when it is so
    obviously a property of the entire number system.
Tortoise: I understand what you mean, but it's not that different from saying “29 is
    prime” or “gold is valuable” – both statements attribute to



Aria with Diverse Variations                                                             407

    a single entity a property which it has only by virtue of being embedded in a
    particular context.
Achilles: I suppose you're right. This "wondrousness" problem is wondrous tricky,
    because of the way in which the numbers oscillate-now increasing, now decreasing.
    The pattern OUGHT to be regular,-yet on the surface it appears to be quite chaotic.
    Therefore, I can well imagine why, as of yet, no one knows of a test for the property
    of wondrousness which is guaranteed to terminate.
Tortoise: Speaking of terminating and nonterminating processes, and those which hover
    in between, I am reminded of a friend of mine, an author, who is at work on a book.
Achilles: Oh, how exciting! What is it called?
Tortoise: Copper, Silver, Gold: an Indestructible Metallic Alloy. Doesn't that sound
    interesting?
Achilles: Frankly, I'm a little confused by the title. After all, what do Copper, Silver, and
    Gold have to do with each other? Tortoise: It seems clear to me.
Achilles: Now if the title were, say, Giraffes, Silver, Gold, or Copper, Elephants, Gold,
    why, I could see it .. .
Tortoise: Perhaps you would prefer Copper, Silver, Baboons?
Achilles: Oh, absolutely! But that original title is a loser. No one would understand it.
Tortoise: I'll tell my friend. He'll be delighted to have a catchier title (as will his
    publisher).
Achilles: I'm glad. But how were you reminded of his book by our discussion?
Tortoise: Ah, yes. You see, in his book there will be a Dialogue in which he wants to
    throw readers off by making them SEARCH for the ending.
Achilles: A funny thing to want to do. How is it done?
Tortoise: You've undoubtedly noticed how some authors go to so much trouble to build
    up great tension a few pages before the end of their stories-but a reader who is
    holding the book physically in his hands can FEEL that the story is about to end.
    Hence, he has some extra information which acts as an advance warning, in a way.
    The tension is a bit spoiled by the physicality of the book. It would be so much better
    if, for instance, there were a lot of padding at the end of novels.
Achilles: Padding?
Tortoise: Yes; what I mean is, a lot of extra printed pages which are not part of the story
    proper, but which serve to conceal the exact location of the end from a cursory
    glance, or from the feel of the book.
Achilles: I see. So a story's true ending might occur, say, fifty or a hundred pages before
    the physical end of the book?
Tortoise: Yes. This would provide an element of surprise, because the reader wouldn't
    know in advance how many pages are padding, and how many are story.
Achilles: If this were standard practice, it might be quite effective. But




Aria with Diverse Variations                                                             408

    there is a problem. Suppose your padding were very obvious-such as a lot of blanks,
    or pages covered with X's or random letters. Then, it would be as good as absent.
Tortoise: Granted. You'd have to make it resemble normal printed pages.
Achilles: But even a cursory glance at a normal page from one story will often suffice to
    distinguish it from another story. So you will have to make the padding resemble the
    genuine story rather closely.
Tortoise: That's quite true. The way I've always envisioned it is this: you bring the story
    to an end; then without any break, you follow it with something which looks like a
    continuation but which is in reality just padding, and which is utterly unrelated to the
    true theme. The padding is, in a way, a "post-ending ending". It may contain
    extraneous literary ideas, having little to do with the original theme.
Achilles: Sneaky! But then the problem is that you won't be able to tell when the real
    ending comes. It'll just blend right into the padding.
Tortoise: That's the conclusion my author friend and I have reached as well. It's a shame,
    for I found the idea rather appealing.
Achilles: Say, I have a suggestion. The transition between genuine story and padding
    material could be made in such a way that, by sufficiently assiduous inspection of
    the text, an intelligent reader will be able to detect where one leaves off and the other
    begins. Perhaps it will take him quite a while. Perhaps there will be no way to
    predict how long it will take ... But the publisher could give a guarantee that a
    sufficiently assiduous search for the true ending will always terminate, even if he
    can't say how long it will be before the test terminates.
Tortoise: Very well-but what does "sufficiently assiduous" mean?
Achilles: It means that the reader must be on the lookout for some small but telltale
    feature in the text which occurs at some point. That would signal the end. And he
    must be ingenious enough to think up, and hunt for, many such features until he
    finds the right one.
Tortoise: Such as a sudden shift of letter frequencies or word lengths? Or a rash of
    grammatical mistakes?
Achilles: Possibly. Or a hidden message of some sort might reveal the true end to a
    sufficiently assiduous reader. Who knows? One could even throw in some
    extraneous characters or events which are inconsistent with the spirit of the
    foregoing story. A naive reader would swallow the whole thing, whereas a
    sophisticated reader would be able to spot the dividing line exactly.
Tortoise: That's a most original idea, Achilles. I'll relay it to my friend, and perhaps he
    can incorporate it in his Dialogue.
Achilles: I would be highly honored.
Tortoise: Well, I am afraid that I myself am growing a little groggy, Achilles. It would be
    well for me to take my leave, while I am still capable of navigating my way home.
Achilles: I am most flattered' that you have stayed up for so long, and at such an odd hour
    of the night, just for my benefit. I assure you that




Aria with Diverse Variations                                                             409

    your number-theoretical entertainment has been a perfect antidote to my usual
    tossing and turning. And who knows-perhaps I may even be able to go to sleep
    tonight. As a token of my gratitude, Mr. T, I would like to present you with a special
    gift.
Tortoise: Oh, don't be silly, Achilles.
Achilles: It is my pleasure, Mr. T. Go over to that dresser; on it, you will see an Asian
    box.

    (The Tortoise moseys over to Achilles' dresser.)

Tortoise. You don't mean this very gold Asian box, do you?
Achilles: That's the one. Please accept it, Mr. T, with my warmest compliments.
Tortoise: Thank you very much indeed, Achilles. Hmm ... Why are all these
    mathematicians' names engraved on the top? What a curious list:

                         De Morgan
                         Abel
                         Boole
                         Br o u w e r
                         Sierpinski
                         Weierstrass

Achilles: I believe it is supposed to be a Complete List of All Great Mathematicians.
    What I haven't been able to figure out is why the letters running down the diagonal
    are so much bolder.
Tortoise: At the bottom it says, "Subtract 1 from the diagonal, to find Bach in Leipzig".
Achilles: I saw that, but I couldn't make head or tail of it. Say, how about a shot of
    excellent whiskey? I happen to have some in that decanter on my shelf.
Tortoise: No, thanks. I'm too tired. I'm just going to head home. (Casually, he opens the
    box.) Say, wait a moment, Achilles-there are one hundred Louis d'or in here!
Achilles: I would be most pleased if you would accept them, Mr. T. Tortoise: But-but
Achilles: No objections, now. The box, the gold-they're yours. And thank you for an
    evening without parallel.
Tortoise: Now whatever has come over you, Achilles? Well, thank you for your
    outstanding generosity and I hope you have sweet dreams about the strange
    Goldbach Conjecture, and its Variation. Good night.

    (And he picks up the very gold Asian box filled with the one hundred Louis d'or, and
    walks towards the door. As he is about to leave, there is a loud knock.)

    Who could be knocking at this ungodly hour, Achilles?
Achilles: I haven't the foggiest idea. It seems suspicious to me. Why don't you go hide
    behind the dresser, in case there's any funny business.




Aria with Diverse Variations                                                          410

Tortoise: Good idea. (Scrambles in behind the dresser.) Achilles: Who's there?
Voice: Open up-it's the cops.
Achilles: Come in, it's open.

    (Two burly policemen walk in, wearing shiny badges.)

Cop: I'm Silva. This is Gould. (Points at his badge.) Is there an Achilles at this address?
Achilles: That's me!
Cop: Well, Achilles, we have reason to believe that there is a very gold Asian box here,
    filled with one hundred Louis d'or. Someone absconded with it from the museum
    this afternoon. Achilles: Heavens to Betsy!
Cop: If it is here, Achilles, since you would be the only possible suspect, I' regret to say
    that I should have to take you into custody. Now I have here a search warrant
Achilles: Oh, sirs, am I ever glad you arrived! All evening long, I have been being
    terrorized by Mr. Tortoise and his very Asian gold box. Now at last you have come
    to liberate me! Please, sirs, just take a look behind that dresser, and there you will
    find the culprit!

    (The cops look behind the dresser and spy the Tortoise huddled behind it, holding
    his very gold Asian box, and trembling.)

Cop: So there it is! And so Mr. Tortoise is the varmint, eh? I never would have suspected
    HIM. But he's caught, red-handed.
Achilles: Haul the villain away, kind sirs! Thank goodness, that's the last I'll have to hear
    of him, and the Very Asian Gold Box!




Aria with Diverse Variations                                                             411


                       BlooP and FlooP and GlooP
                             Self-Awareness and Chaos
BLOOP, FLOOP, AND GLOOP are not trolls, talking ducks, or the sounds made by a
sinking ship-they are three computer languages, each one with is own special purpose.
These languages were invented specially for this chapter. They will be of use in
explaining some new senses of the word 'recursive -in particular, the notions of primitive
recursivity and general recursivity. They will prove very helpful in clarifying the
machinery of self-reference in TNT.
         We seem to be making a rather abrupt transition from brains and hinds to
technicalities of mathematics and computer science. Though the transition is abrupt in
some ways, it makes some sense. We just saw how a certain kind of self-awareness
seems to be at the crux of consciousness. Vow we are going to scrutinize "self-
awareness" in more formal settings, such as TNT. The gulf between TNT and a mind is
wide, but some of the ideas will be most illuminating, and perhaps metaphorically
transportable back to our thoughts about consciousness.
         One of the amazing things about TNT's self-awareness is that it is intimately
connected to questions about order versus chaos among the natural numbers. In
particular, we shall see that an orderly system of sufficient complexity that it can mirror
itself cannot be totally orderly-it must contain some strange, chaotic features. For readers
who have some Achilles in them, this will be hard to take. However, there is a "magical"
compensation: there is a kind of order to the disorder, which is now its own field of
study, called "recursive function theory". Unfortunately, we will not be able to do much
more than hint at the fascination of this subject.

                       Representability and Refrigerators
Phrases such as "sufficiently complex", "sufficiently powerful" and the like lave cropped
up quite often earlier. Just what do they mean? Let us go back to the battle of the Crab
and Tortoise, and ask, "What qualifies something as a record player?" The Crab might
claim that his refrigerator s a "Perfect" record player. Then to prove it, he could set any
record whatsoever atop it, and say, "You see-it's playing it!" The Tortoise, if he wanted to
counter this Zen-like act, would have to reply, "No-your refrigerator is too low-fidelity to
be counted as a phonograph: it cannot reproduce sounds-at all (let alone its self-breaking
sound)." The Tortoise




BlooP and FlooP and GlooP                                                               413

can only make a record called "I Cannot Be Played on Record Player X" provided that
Record Player X is really a record player! The Tortoise's method is quite insidious, as it
plays on the strength, rather than on the weakness, of the system. And therefore he
requires "sufficiently hi-fi" record players.
        Ditto for formal versions of number theory. The reason that TNT is a
formalization of N is that its symbols act the right way: that is, its theorems are not silent
like a refrigerator-they speak actual truths of N. Of course, so do the theorems of the pq-
system. Does it, too, count as "a formalization of number theory", or is it more like a
refrigerator? Well, it is a little better than a refrigerator, but it is still pretty weak. The pq-
system does not include enough of the core truths of N to count as "a number theory".
        What, then, are these "core truths" of N? They are the primitive recursive truths;
that means they involve only predictably terminating calculations. These core truths
serve for N as Euclid's first four postulates served for geometry: they allow you to throw
out certain candidates before the game begins, on the grounds of "insufficient power".
From here on out, the representability of all primitive recursive truths will be the
criterion for calling a system "sufficiently powerful".

                           Ganto's Ax in Metamathematics
The significance of the notion is shown by the following key fact: If you have a
sufficiently powerful formalization of number theory, then Gödel’s method is applicable,
and consequently your system is incomplete. If, on the other hand, your system is not
sufficiently powerful (i.e., not all primitive recursive truths are theorems), then your
system is, precisely by virtue of that lack, incomplete. Here we have a reformulation of
"Ganto's Ax" in metamathematics: whatever the system does, Gödel’s Ax will chop its
head off! Notice also how this completely parallels the high-fidelity-versus-low fidelity
battle in the Contracrostipunctus.
Actually, it turns out that much weaker systems are still vulnerable to the Gödel method;
the criterion that all primitive recursive truths need be represented as theorems is far too
stringent. It is a little like a thief who will only rob "sufficiently rich" people, and whose
criterion is that the potential victim should be carrying at least a million dollars in cash.
In the case of TNT, luckily, we will be able to act in our capacity as thieves, for the
million in cash is there-which is to say, TNT does indeed contain all primitive recursive
truths as theorems.
        Now before we plunge into a detailed discussion of primitive recursive functions
and predicates, I would like to tie thee themes of this Chapter to themes from earlier
Chapters, so as to provide a bit better motivation.

                   Finding Order by Choosing the Right Filter
We saw at a very early stage that formal systems can be difficult and unruly beasts
because they have lengthening and shortening rules, which can




BlooP and FlooP and GlooP                                                                     414

possibly lead to never-ending searches among strings. The discovery of Gödel-numbering
showed that any search for a string having a special typographical property has an
arithmetical cousin: an isomorphic search for an integer with a corresponding special
arithmetical property. Consequently, the quest for decision procedures for formal systems
involves solving the mystery of unpredictably long searches- chaos-among the integers.
Now in the Aria with Diverse Variations, I gave perhaps too much weight to apparent
manifestations of chaos in problems about integers. As a matter of fact, people have
tamed wilder examples of apparent chaos than the "wondrousness" problem, finding them
to be quite gentle beasts after all. Achilles' powerful faith in the regularity and
predictability of numbers should therefore be accorded quite a bit of respect-especially as
it reflects the beliefs of nearly all mathematicians up till the 1930's. To show why order
versus chaos is such a subtle and significant issue, and to tie it in with questions about the
location and revelation of meaning, I would like to quote a beautiful and memorable
passage from Are Quanta Real?-a Galilean Dialogue by the late J. M. Jauch:

        SALVIATI Suppose I give you two sequences of numbers, such as
             78539816339744830961566084...
  And

               1, -1/3, +1/5, -1/7, +1/9, -1/11, +1/13, -1/15, ...

      If I asked you, Simplicio, what the next number of the first sequence is, what
  would you say?
      SIMPLICIO I could not tell you. I think it is a random sequence and that there is
  no law in it.
      SALVIATI And for the second sequence?
      SIMPLICIO That would be easy. It must be +1/17.
      SALVIATI Right. But what would you say if I told you that the first
      sequence is also constructed by a law and this law is in fact identical with the
      one you have just discovered for the second sequence? SIMPLICIO This does not
  seem probable to me.
      SALVIATI But it is indeed so, since the first sequence is simply the beginning of
  the decimal fraction [expansion] of the sum of the second. Its value is Tr/4.
      SIMPLICIO You are full of such mathematical tricks, but I do not see what this
  has to do with abstraction and reality.
      SALVIATI The relationship with abstraction is easy to see. The first sequence
  looks random unless one has developed through a process of abstraction a kind of
  filter which sees a simple structure behind the apparent randomness.
      It is exactly in this manner that laws of nature are discovered. Nature
      presents us with a host of phenomena which appear mostly as chaotic randomness
  until we select some significant events, and abstract from their particular, irrelevant
  circumstances so that they become idealized. Only then can they exhibit their true
  structure in full splendor.
      SAGREDO This is a marvelous idea! It suggests that when we try to understand
  nature, we should look at the phenomena as if they were messages to be




BlooP and FlooP and GlooP                                                                 415

  understood. Except that each message appears to be random until we establish a code
  to read it. This code takes the form of an abstraction, that is, we choose to ignore
  certain things as irrelevant and we thus partially select the content of the message by
  a free choice. These irrelevant signals form the "background noise," which will limit
  the accuracy of our message.
     But since the code is not absolute there may be several messages in the same raw
  material of the data, so changing the code will result in a message of' equally deep
  significance in something that was merely noise before, and conversely: In a new
  code a former message may be devoid of meaning.
     Thus a code presupposes a free choice among different, complementary aspects,
  each of which has equal claim to reality, if I may use this dubious word.
     Some of these aspects may be completely unknown to us now but they may
  reveal themselves to an observer with a different system of abstractions.
     But tell me, Salviati, how can we then still claim that we discover something out
  there in the objective real world? Does this not mean that we are merely creating
  things according to our own images and that reality is only within ourselves?
     SALVIATI I don't think that this is necessarily so, but it is a question which
  requires deeper reflection.'

Jauch is here dealing with messages that come not from a "sentient being" but from
nature itself. The questions that we raised in Chapter VI on the relation of meaning to
messages can be raised equally well with messages from nature. Is nature chaotic, or is
nature patterned? And what is the role of intelligence in determining the answer to this
question?
        To back off from the philosophy, however, we can consider the point about the
deep regularity of an apparently random sequence. Might the function Q(n) from Chapter
V have a simple, nonrecursive explanation, too? Can every problem, like an orchard, be
seen from such an angle that its secret is revealed? Or are there some problems in number
theory which, no matter what angle they are seen from, remain mysteries?
        With this prologue, I feel it is time to move ahead to define the precise meaning
of the term "predictably long search". This will be accomplished in terms of the language
B1ooP.

                   Primordial Steps of the Language BlooP
Our topic will be searches for natural numbers which have various properties. In order to
talk about the length of any search, we shall have to define some primordial steps, out of
which all searches are built, so that length can be measured in terms of number of steps.
Some steps which we might consider primordial are:

    adding any two natural numbers;
    multiplying any two natural numbers;
    determining if two numbers are equal;
    determining the larger (smaller) of two numbers.




BlooP and FlooP and GlooP                                                              416

                              Loops and Upper Bounds
If we try to formulate a test for, say, primality in terms of such steps, we shall soon see
that we have to include a control structure-that is, descriptions of the order to do things
in, when to branch back and try something again, when to skip over a set of steps, when
to stop, and similar matters.
It is typical of any algorithm-that is, a specific delineation of how to carry out a task-that
it includes a mixture of (1) specific operations to be performed, and (2) control
statements. Therefore, as we develop our language for expressing predictably long
calculations, we shall have to incorporate primordial control structures also. In fact, the
hallmark of BlooP is its limited set of control structures. It does not allow you to branch
to arbitrary steps, or to repeat groups of steps without limit; in BlooP, essentially the only
control structure is the bounded loop: a set of instructions which can be executed over
and over again, up to a predefined maximum number of times, called the upper bound, or
ceiling, of the loop. If the ceiling were 300, then the loop might be executed 0, 7, or 300
times-but not 301.
Now the exact values of all the upper bounds in a program need not be put in numerically
by the programmer-indeed, they may not be known in advance. Instead, any upper bound
may be determined by calculations carried out before its loop is entered. For instance, if
you wanted to calculate the value of 2"', there would be two loops. First, you evaluate 3",
which involves n multiplications. Then, you put 2 to that power, which involves 3"
multiplications. Thus, the upper bound for the second loop is the result of the calculation
of the first loop.
       Here is how you would express this in a BlooP program:

      DEFINE PROCEDURE "TWO-TO-THE-THREE-TO-THE" [N]:
      BLOCK 0: BEGIN
                 CELL(O) <= 1;
                 LOOP N TIMES:
                 BLOCK 1: BEGIN
                               CELL(0) ' 3 x CELL(O);
                 BLOCK 1: END;
                 CELL(1) <= 1;
                 LOOP CELL(O) TIMES:
                 BLOCK 2: BEGIN
                               CELL(1) \# 2 X CELL(l );
                BLOCK 2: END;
                OUTPUT <= CELL( I );
       BLOCK 0: END.

                                Conventions of BlooP
Now it is an acquired skill to be able to look at an algorithm written in a computer
language, and figure out what it is doing. However, I hope that this algorithm is simple
enough that it makes sense without too much



BlooP and FlooP and GlooP                                                                 417

scrutiny. A procedure is defined, having one input parameter, N; its output is the desired
value.
         This procedure definition has what is called block structure, which means that
certain portions of it are to be considered as units, or blocks. All the statements in a block
get executed as a unit. Each block has a number (the outermost being BLOCK 0), and is
delimited by a BEGIN and an END. In our example, BLOCK 1 and BLOCK 2 contain
just one statement each but shortly you will see longer blocks. A LOOP statement always
means to execute the block immediately under it repeatedly. As can be seen above,
blocks can be nested.
         The strategy of the above algorithm is as described earlier. You begin by taking
an auxiliary variable, called CELL(O); you set it initially to 1, and then, in a loop, you
multiply it repeatedly by 3 until you've done so exactly N times. Next, you do the
analogous thing for CELL(1)-set it to 1, multiply by 2 exactly CELL(O) times, then
quit. Finally, you set OUTPUT to the value of CELL(1). This is the value returned to the
outside world-the only externally visible behavior of the procedure.
         A number of points about the notation should be made here. First, the meaning of
the left-arrow <= is this:

Evaluate the expression to its right, then take the result and set the CELL (or OUTPUT)
           on its left to that value.

So the meaning of a command such as CELL(1) <= 3 X CELL(1) is to triple the value
stored in CELL(1). You may think of each CELL as being a separate word in the
memory of some computer. The only difference between a CELL and a true word is that
the latter can only hold integers up to some finite limit, whereas we allow a CELL to
hold any natural number, no matter how big.
        Every procedure in BlooP, when called, yields a value-namely the value of the
variable called OUTPUT. At the beginning of execution of any procedure, it is assumed
as a default option that OUTPUT has the value 0. That way, even if the procedure never
resets OUTPUT at all, OUTPUT has a well-defined value at all times.

                           IF-Statements and Branching
Now let us look at another procedure which will show us some other features of BlooP
which give it more generality. How do you find out, knowing only how to add, what the
value of M - N is? The trick is to add various numbers onto N until you find the one
which yields M. However, what happens if M is smaller than N? What if we are trying to
take 5 from 2? In the domain of natural numbers, there is no answer. But we would like
our B1ooP procedure to give an answer anyway-let's say 0. Here, then, is a BlooP
procedure which does subtraction:




BlooP and FlooP and GlooP                                                                 418

DEFINE PROCEDURE "MINUS" [M,N]:
BLOCK 0: BEGIN
     IF M < N, THEN:
     QUIT BLOCK 0;
     LOOP AT MOST M + 1 TIMES:
     BLOCK 1: BEGIN
                 IF OUTPUT + N = M, THEN:
                 ABORT LOOP 1;
                 OUTPUT, <= OUTPUT + 1;
     BLOCK 1: END;
BLOCK 0: END.

        Here we are making use of the implicit feature that OUTPUT begins at 0. If M is
less than N, then the subtraction is impossible, and we simply jump to the bottom of
BLOCK 0 right away, and the answer is 0. That is what is meant by the line QUIT
BLOCK 0. But if M is not less than N, then we skip over that QUIT-statement, and
carry out the next command in sequence (here, a LOOP-statement). That is how IF-
statements always work in BlooP.
        So we enter LOOP 1, so called because the block which it tells us to repeat is
BLOCK 1. We try adding 0 to N, then 1, 2, etc., until we find a number that gives M. At
that point, we ABORT the loop we are in, meaning we jump to the statement
immediately following the END which marks the bottom of the loop's block. In this case,
that jump brings us just below BLOCK 1: END, which is to say, to the last statement of
the algorithm, and we are done. OUTPUT now contains the correct answer.
        Notice that there are two distinct instructions for jumping downwards: QUIT, and
ABORT. The former pertains to blocks, the latter to loops. QUIT BLOCK n means to
jump to the last line of BLOCK n, whereas ABORT LOOP n means to jump just below
the last line of BLOCK n. This distinction only matters when you are inside a loop and
want to continue looping but to quit the block this time around. Then you can say QUIT
and the proper thing will happen.
        Also notice that the words AT MOST now precede the upper bound of the loop,
which is a warning that the loop may be aborted before the upper bound is reached.

                                Automatic Chunking
Now there are two last features of BlooP to explain, both of them very important. The
first is that, once a procedure has been defined, it may be called inside later procedure
definitions. The effect of this is that once an operation has been defined in a procedure, it
is considered as simple as a primordial step. Thus, BlooP features automatic chunking.
You might compare it to the way a good ice skater acquires new motions: not by defining
them as long sequences of primordial muscle-actions, but in terms of previously learned
motions, which were themselves learned as compounds of earlier




BlooP and FlooP and GlooP                                                                419

learned motions, etc.-and the nestedness, or chunkedness, can go back many layers until
you hit primordial muscle-actions And thus, the repertoire of BlooP programs, like the
repertoire of a skater's tricks, grows, quite literally, by loops and bounds.

                                      BlooP Tests
The other feature of BlooP is that certain procedures can have YES or NO as their
output, instead of an integer value. Such procedures are tests, rather than functions. To
indicate the difference, the name of a test must terminate in a question mark. Also, in a
test, the default option for OUTPUT is not 0, of course, but NO.
Let us see an example of these last two features of BlooP in an algorithm which tests its
argument for primality:

       DEFINE PROCEDURE "PRIME?" [N]:
       BLOCK 0: BEGIN
                  IF N = 0, THEN:
                  QUIT BLOCK 0;
                  CELL(0) <= 2;
                  LOOP AT MOST MINUS [N,2] TIMES:
                  BLOCK 1: BEGIN
                               IF REMAINDER [N,CELL(O)] = 0, THEN:
                               QUIT BLOCK 0;
                               CELL(O) <= CELL(O) + 1;
                  BLOCK 1: END;
                  OUTPUT <= YES;
       BLOCK 0: END.

Notice that I have called two procedures inside this algorithm: MINUS and
REMAINDER. (The latter is presumed to have been previously defined, and you may
work out its definition yourself.) Now this test for primality works by trying out potential
factors of N one by one, starting at 2 and increasing to a maximum of N - 1. In case any
of them divides N exactly (i.e., gives remainder 0), then we jump down to the bottom,
and since OUTPUT still has its default value at this stage, the answer is NO. Only if N
has no exact divisors will it survive the entirety of LOOP 1; then we will emerge
smoothly at the statement OUTPUT <= YES, which will get executed, and then the
procedure is over.

               BlooP Programs Contain Chains of Procedures
We have seen how to define procedures in BlooP; however, a procedure definition is only
a part of a program. A program consists of a chain of procedure definitions (each only
calling previously defined procedures), optionally followed by one or more calls on the
procedures defined. Thus, an




BlooP and FlooP and GlooP                                                               420

example of a full B1ooP program would be the definition of the procedure TWO-TO-
THE-THREE-TO-THE, followed by the call

                         TWO-TO-THE-THREE-TO-THE [2]

which would yield an answer of 512.
         If you have only a chain of procedure definitions, then nothing ever gets executed;
they are all just waiting for some call, with specific numerical values, to set them in
motion. It is like a meat grinder waiting for some meat to grind-or rather, a chain of meat
grinders all linked together, each of which is fed from earlier ones ... In the case of meat
grinders, the image is perhaps not so savory; however, in the case of BlooP programs,
such a construct is quite important, and we will call it a "call-less program". This notion
is illustrated in Figure 72.
         Now B1ooP is our language for defining predictably terminating calculations. The
standard name for functions which are B1ooP-computable is primitive recursive
functions; and the standard name for properties which can be detected by B1ooP-tests is
primitive recursive predicates. Thus, the function 23n is a primitive recursive function;
and the statement "n is a prime number" is a primitive recursive predicate.
         It is clear intuitively that the Goldbach property is primitive recursive, and to
make that quite explicit, here is a procedure definition in BlooP, showing how to test for
its presence or absence:

DEFINE PROCEDURE "GOLDBACH?" [N]:
BLOCK 0: BEGIN
           CELL(0)    2;
           LOOP AT MOST N TIMES:
           BLOCK 1: BEGIN
                      IF {PRIME? [CELL(O)]
                         AND PRIME? [MINUS [N,CELL(0)]]},
                            THEN:
                            BLOCK 2: BEGIN
                                  OUTPUT,\# YES;
                                  QUIT BLOCK 0-,
                            BLOCK 2: END
                            CELL(0) <= CELL(0) +
           BLOCK 1: END;
BLOCK 0: END.


As usual, we assume NO until proven YES, and we do a brute force search among pairs
of numbers which sum up to N. If both are prime, we quit the outermost block; otherwise
we just go back and try again, until all possibilities are exhausted.
(Warning: The fact that the Goldbach property is primitive recursive does not make the
question “Do all numbers have the Goldbach property?” a simple question—far from it!)




BlooP and FlooP and GlooP                                                               421

FIGURE 72. The structure of a call-less BlooP program. For this program to be self-
contained, each procedure definition may only call procedures defined above it.

                                 Suggested Exercises
Can you write a similar B1ooP procedure which tests for the presence or absence of the
Tortoise property (or the Achilles property)? If so, do it. If not, is it merely because you
are ignorant about upper bounds, or could it be that there is a fundamental obstacle
preventing the formulation of such an algorithm in BlooP? And what about the same
questions, with respect to the property of wondrousness, defined in the Dialogue?
Below, I list some functions and properties, and you ought to take the time to determine
whether you believe they are primitive recursive (BlooP-programmable) or not. This
means that you must carefully consider what kinds of operations will be involved in the
calculations which they require, and whether ceilings can be given for all the loops
involved.

FACTORIAL [N] = NI (the factorial of N)
 (e.g., FACTORIAL [4] = 24)
REMAINDER [M,N] = the remainder upon dividing M by N
 (e.g., REMAINDER [24,7] = 3)
PI-DIGIT [N] = the Nth digit of pi, after the decimal point
 (e.g. PI-DIGIT [1] = 1,
        PI-DIGIT [2] = 4
        PI-DIGIT [1000000] = 1



BlooP and FlooP and GlooP                                                               422

FIBO [N] = the Nth Fibonacci number
(e.g., FIBO [9] = 34)
PRIME-BEYOND [N[ = the lowest prime beyond N
(e.g., PRIME-BEYOND [33] = 37)
PERFECT [N] = the Nth "perfect" number (a number such as 28 whose divisors sum up
to itself: 28 = 1 + 2 + 4 + 7 + 14)
(e.g., PERFECT [2] = 28)
PRIME? [N] = YES if N is prime, otherwise NO.
PERFECT? [N] = YES if N is perfect, otherwise NO.
TRIVIAL? [A,B,C,N] = YES if A"+B" = Cn is correct; otherwise NO.
(e.g., TRIVIAL? [3,4,5,2] = YES,
         TRIVIAL? [3,4,5,3] = NO)
PIERRE? [A,B,C] = YES if A"+B" = C" is satisfiable for some value of N greater than
1, otherwise NO.
(e.g., PIERRE? [3,4,5] = YES,
         PIERRE? [1,2,3] = NO)
FERMAT? [N] = YES if A"+B" = CN is satisfied by some positive
values of A, B, C; otherwise NO.
(e.g., FERMAT? [2] = YES)
TORTOISE-PAIR? [M,N] = YES if both M and M + N are prime, otherwise NO.
(e.g., ORTOISE-PAIR [5,1742] = YES,
         TORTOISE-PAIR [5,100] = NO)
TORTOISE? [N] = YES if N is the difference of two primes, otherwise NO.
(e.g., TORTOISE [1742] = YES,
         TORTOISE [7] = NO)
MIU-WELL-FORMED? [N] = YES if N, when seen as a string of the MIU-System, is
well-formed; otherwise NO.
(e.g., MIU-WELL-FORMED? [310] = YES,
         MIU-WELL-FORMED? [415] = NO)
MIU-PROOF-PAIR? [M,N] = YES If M, as seen as a sequence of strings of the MIU-
system, is a derivation of N, as seen as a string of the MIU-system; otherwise NO.
(e.g., MIU-PROOF-PAIR? [3131131111301,301] = YES,
         MIU-PROOF-PAIR? [311130,30] = NO)
MIU-THEOREM? [N] = YES if N, seen as a MIU-system string, is a theorem;
otherwise NO.
(e.g., MIU-THEOREM? [311] = YES,
         MIU-THEOREM? [30] = NO,
         MIU-THEOREM? [701] = NO)
TNT-THEOREM? [N] = YES if N, seen as a TNT-string, is a theorem.
(e.g., TNT-THEOREM? [666111666] = YES,
         TNT-THEOREM? [123666111666] = NO,
         TNT-THEOREM? [7014] = NO)




BlooP and FlooP and GlooP                                                      423

FALSE? [N] = YES if N, seen as a TNT-string, is a false statement of number theory;
otherwise NO.
(e.g., FALSE? [6661 1 1666] = NO,
       FALSE? [2236661 1 1666] = YES,
       FALSE? [7014] = NO)

The last seven examples are particularly relevant to our future metamathematical
explorations, so they highly merit your scrutiny.

                          Expressibility and Representability

Now before we go on to some interesting questions about BlooP and are led to its
relative, FlooP, let us return to the reason for introducing BlooP in the first place, and
connect it to TNT. Earlier, I stated that the critical mass for Gödel’s method to be
applicable to a formal system is attained when all primitive recursive notions are
representable in that system. Exactly what does this mean? First of all, we must
distinguish between the notions of representability and expressibility. Expressing a
predicate is a mere matter of translation from English into a strict formalism. It has
nothing to do with theoremhood. For a predicate to be represented, on the other hand, is a
much stronger notion. It means that

     (1) All true instances of the predicate are theorems;
     (2) All false instances are nontheorems.

By "instance", I mean the string produced when you replace all free variables by
numerals. For example, the predicate m + n = k is represented in the pq-system, because
each true instance of the predicate is a theorem, each false instance is a nontheorem. Thus
any specific addition, whether true or false, translates into a decidable string of the pq-
system. However, the pq-system is unable to express-let alone represent-any other
properties of natural numbers. Therefore it would be a weak candidate indeed in a
competition of systems which can do number theory.
        Now TNT has the virtue of being able to express virtually any number-theoretical
predicate; for example, it is easy to write a TNT-string which expresses the predicate "b
has the Tortoise property". Thus, in terms of expressive power, TNT is all we want.
        However, the question "Which properties are represented in TNT?" is Precisely
the question "How powerful an axiomatic system is TNT?" Are all Possible predicates
represented in TNT? If so, then TNT can answer any question of number theory; it is
complete.

          Primitive Recursive Predicates Are Represented in TNT

Now although completeness will turn out to be a chimera. TNT is at least complete with
respect to primitive recursive predicates. In other words, any statement of number theory
whose truth or falsity can be decided by a




BlooP and FlooP and GlooP                                                              424

computer within a predictable length of time is also decidable inside TNT. Or, one final
restatement of the same thing:

      If a BlooP test can be written for some property of natural numbers, then that
      property is represented in TNT.

         Are There Functions Which Are Not Primitive Recursive?
Now the kinds of properties which can be detected by BlooP tests are widely varied,
including whether a number is prime or perfect, has the Goldbach property, is a power of
2, and so on and so forth. It would not be crazy to wonder whether every property of
numbers can be detected by some suitable BlooP program. The fact that, as of the present
moment, we have no way of testing whether a number is wondrous or not need not
disturb us too much, for it might merely mean that we are ignorant about wondrousness,
and that with more digging around, we could discover a universal formula for the upper
bound to the loop involved. Then a BlooP test for wondrousness could be written on the
spot. Similar remarks could be made about the Tortoise property.
        So the question really is, "Can upper bounds always be given for the length of
calculations-or, is there an inherent kind of jumbliness to the natural number system,
which sometimes prevents calculation lengths from being predictable in advance?" The
striking thing is that the latter is the case, and we are about to see why. It is the sort of
thing that would have driven Pythagoras, who first proved that the square root of 2 is
irrational, out of his mind. In our demonstration, we will use the celebrated diagonal
method, discovered by Georg Cantor, the founder of set theory.

                  Pool B, Index Numbers, and Blue Programs
We shall begin by imagining a curious notion: the pool of all possible BlooP programs.
Needless to say, this pool-"Pool B"-is an infinite one. We want to consider a subpool of
Pool B, obtained by three successive filtering operations. The first filter will retain for us
only call-less programs. From this subpool we then eliminate all tests, leaving only
functions. (By the way, in call-less programs, the last procedure in the chain determines
whether the program as a whole is considered a test, or a function.) The third filter will
retain only functions which have exactly one input parameter. (Again referring to the
final procedure in the chain.) What is left?

       A complete pool of all call-less BlooP programs which calculate functions of
       exactly one input parameter.

Let us call these special BlooP programs Blue Programs.
         What we would like to do now is to assign an unambiguous index
number to each Blue Program. How can this be done? The easiest way—we shall use it—
is to list them in order of length: the shortest possible. Blue




BlooP and FlooP and GlooP                                                                 425

Program being \# 1, the second shortest being \#2, etc. Of course, there will be many
programs tied for each length. To break such ties, we use alphabetical order. Here,
"alphabetical order" is taken in an extended sense, where the alphabet includes all the
special characters of BlooP, in some arbitrary order, such as the following:

            A B C D E F G H I J K L M N
            O P Q R S T U V W X Y Z + x
            0 1 2 3 4 5 6 7 8 9 <= = < >
            ( ) [ ] { } - ´ ? : ; , .

-and at the end comes the lowly blank! Altogether, fifty-six characters. For convenience's
sake, we can put all Blue Programs of length 1 in Volume 1, programs of 2 characters in
Volume 2, etc. Needless to say, the first few volumes will be totally empty, while later
volumes will have many, many entries (though each volume will only have a finite
number). The very first Blue Program would be this one:

            DEFINE PROCEDURE "A" [B]:
            BLOCK 0: BEGIN
            BLOCK 0: END.

This rather silly meat grinder returns a value of 0 no matter what its input is. It occurs in
Volume 56, since it has 56 characters (counting necessary blanks, including blanks
separating successive lines).
         Soon after Volume 56, the volumes will get extremely fat, because there are just
so many millions of ways of combining symbols to make Blue BlooP programs. But no
matter-we are not going to try to print out this infinite catalogue. All that we care about is
that, in the abstract, it is well-defined, and that each Blue BlooP program therefore has a
unique and definite index number. This is the crucial idea.
         Let us designate the function calculated by the kth Blue Program this way:

                                   Blueprogram{\# k} [N]

Here, k is the index number of the program, and N is the single input parameter. For
instance, Blue Program \# 12 might return a value twice the size of its input:

                              Blueprogram{\#12} [N] = 2 x N

The meaning of the equation above is that the program named on the left-hand side
returns the same value as a human would calculate from the ordinary algebraic
expression on the right-hand side. As another example, perhaps the 5000th Blue Program
calculates the cube of its input parameter:

                              Blueprogram{\#5000} [N] = N3




BlooP and FlooP and GlooP                                                                 426

                                The Diagonal Method
Very well-now we apply the "twist": Cantor's diagonal method. We shall take this
catalogue of Blue Programs and use it to define a new function of one variable-Bluediag
[N]-which will turn out not to be anywhere in the list (which is why its name is in italics).
Yet Bluediag will clearly be a well-defined, calculable function of one variable, and so
we will have to conclude that functions exist which simply are not programmable in
BlooP.
      Here is the definition of Bluediag ~N]:

      Equation (1) ... Bluediag [N] = 1 + Blueprogram{\#N} [N]

The strategy is: feed each meat grinder with its own index number, then add 1 to the
output. To illustrate, let us find Bluediag [12]. We saw that Blueprogram{\# 12} is the
function 2N; therefore, Bluediag [12] must have the value 1 + 2 x 12, or 25. Likewise,
Bluediag [5000] would have the value 125,000,000,001, since that is 1 more than the
cube of 5000. Similarly, you can find Bluediag of any particular argument you wish.
        The peculiar thing about Bluediag [N] is that it is not represented in the catalogue
of Blue Programs. It cannot be. The reason is this. To be a Blue Program, it would have
to have an index number-say it were Blue Program \# X. This assumption is expressed by
writing

       Equation (2) ...        Bluediag [N] = Blueprogram{\# X} [N]

But there is an inconsistency between the equations (1) and (2). It becomes apparent at
the moment we try to calculate the value of Bluediag [ X], for we can do so by letting N
take the value of X in either of the two equations. If we substitute into equation (1), we
get:

       Bluediag [ X] = 1 + Blueprogram{\# X} [ X]

But if we substitute into equation (2) instead, we get:

       Bluediag [ X] = Blueprogram{\# X} [ X]

Now Bluediag [ X] cannot be equal to a number and also to the successor of that number.
But that is what the two equations say. So we will have to go back and erase some
assumption on which the inconsistency is based. The only possible candidate for erasure
is the assumption expressed by Equation (2): that the function Bluediag [N] is able to be
coded up as a Blue BlooP program. And that is the proof that Bluediag lies outside the
realm of primitive recursive functions. Thus, we have achieved our aim of destroying
Achilles' cherished but naive notion that every number-theoretical function must be
calculable within a predictable number of steps.
        There are some subtle things going on here. You might ponder this, for instance:
the number of steps involved in the calculation of Bluediag [N],for each specific value of
N, is predictable—but the different methods of prediction cannot all be united into a
general recipe for predict


BlooP and FlooP and GlooP                                                                427

ing the length of calculation of Bluediag [N]. This is an "infinite conspiracy", related to
the Tortoise's notion of "infinite coincidences", and also to w-incompleteness. But we
shall not trace out the relations in detail.

                      Cantor's Original Diagonal Argument
Why is this called a diagonal argument? The terminology comes from Cantor's original
diagonal argument, upon which many other arguments (such as ours) have subsequently
been based. To explain Cantor's original argument will take us a little off course, but it is
worthwhile to do so. Cantor, too, was concerned with showing that some item is not in a
certain list. Specifically, what Cantor wanted to show was that if a "directory" of real
numbers were made, it would inevitably leave some real numbers out-so that actually, the
notion of a complete directory of real numbers is a contradiction in terms.
        It must be understood that this pertains not just to directories of finite size, but
also to directories of infinite size. It is a much deeper result than the statement "the
number of reals is infinite, so of course they cannot be listed in a finite directory". The
essence of Cantor's result is that there are (at least) two distinct types of infinity: one kind
of infinity describes how many entries there can be in an infinite directory or table, and
another describes how many real numbers there are (i.e., how many points there are on a
line, or line segment)-and this latter is "bigger", in the sense that the real numbers cannot
be squeezed into a table whose length is described by the former kind of infinity. So let
us see how Cantor's argument involves the notion of diagonal, in a literal sense.
        Let us consider just real numbers between 0 and 1. Assume, for the sake of
argument, that an infinite list could be given, in which each positive integer N is matched
up with a real number r(N) between 0 and 1, and in which each real number between 0
and 1 occurs somewhere down the line. Since real numbers are given by infinite
decimals, we can imagine that the beginning of the table might look as follows:

r(1):   .1     4       1       5       9       2       6       5       3
r(2):   .3     3       3       3       3       3       3       3       3
r(3):   .7     1       8       2       8       1       8       2       8
r(4):   .4     1       4       2       1       3       5       6       2
r(5):   .5     0       0       0       0       0       0       0       0

The digits that run down the diagonal are in boldface: 1, 3, 8, 2, 0.... Now those diagonal
digits are going to be used in making a special real number d, which is between 0 and 1
but which, we will see, is not in the list. To make d, you take the diagonal digits in order,
and change each one of them to some other digit. When you prefix this sequence of digits
by a decimal point you have d. There are of course many ways of changing a digit to
some other digit, and correspondingly many different d´s. Suppose for,




BlooP and FlooP and GlooP                                                                  428

example, that we subtract 1 from the diagonal digits (with the convention that 1 taken
from 0 is 9). Then our number d will be:

                                        .0 2 7 1 9 .

Now, because of the way we constructed it,

            d's 1st digit is not the same as the 1st digit of r(1);
            d's 2nd digit is not the same as the 2nd digit of r(2);
            d's 3rd digit is not the same as the 3rd digit of r(3);

            ... and so on.

Hence,

            d is different from r(1);
            d is different from r(2);
            d is different from r(3);
            ... and soon.

In other words, d is not in the list!

                     What Does a Diagonal Argument Prove?
Now comes the crucial difference between Cantor's proof and our proofit is in the matter
of what assumption to go back and undo. In Cantor's argument, the shaky assumption
was that such a table could be drawn up. Therefore, the conclusion warranted by the
construction of d is that no exhaustive table of reals can be drawn up after all-which
amounts to saying that the set of integers is just not big enough to index the set of reals.
On the other hand, in our proof, we know that the directory of Blue BlooP programs can
be drawn up-the set of integers is big enough to index the set of Blue BlooP programs.
So, we have to go back and retract some shakier idea which we used. And that idea is that
Bluediag [N] is calculable by some program in BlooP. This is a subtle difference in the
application of the diagonal method.
        It may become clearer if we apply it to the alleged "List of All Great
Mathematicians" in the Dialogue-a more concrete example. The diagonal itself is
"Dboups". If we perform the desired diagonal-subtraction, we will get "Cantor". Now two
conclusions are possible. If you have an unshakable belief that the list is complete, then
you must conclude that Cantor is not a Great Mathematician, for his name differs from all
those on the list. On the other hand, if you have an unshakable belief that Cantor is a
Great Mathematician, then you must conclude that the List of All Great Mathematicians
is incomplete, for Cantor's name is not on the list! (Woe to those who have unshakable
beliefs on both sides!) The former case corresponds to our proof that Bluediag [N] is not
primitive recursive; the latter case corresponds to Cantor´s proof that the list of reals is
incomplete;




BlooP and FlooP and GlooP                                                               429

                                FIGURE 73. Georg Cantor

Cantor's proof uses a diagonal in the literal sense of the word. Other "diagonal" proofs are
based on a more general notion, which is abstracted from the geometric sense of the
word. The essence of the diagonal method is the fact of using one integer in two different
ways-or, one could say, using one integer on two different levels-thanks to which one can
construct an item which is outside of some predetermined list. One time, the integer
serves as a vertical index, the other time as a horizontal index. In Cantor's construction
this is very clear. As for the function Bluediag [N], it involves using one integer on two
different levels-first, as a Blue Program index number; and second, as an input parameter.

           The Insidious Repeatability of the Diagonal Argument
At first, the Cantor argument may seem less than fully convincing. Isn't there some way
to get around it? Perhaps by throwing in the diagonally constructed number d, one might
obtain an exhaustive list. If you consider this idea, you will see it helps not a bit to throw
in the number d, for as soon as you assign it a specific place in the table, the diagonal
method becomes applicable to the new table, and a new missing number d' can be
constructed, which is not in the new table. No matter how many times you repeat the
operation of constructing a number by the diagonal method and then throwing it in to
make a "more complete" table, you still are caught on the ineradicable hook of Cantor’s
method. You might even try to build a table of reals which tries to outwit the Cantor
diagonal method by taking




BlooP and FlooP and GlooP                                                                 430

the whole trick, lock, stock, and barrel, including its insidious repeatability, into account
somehow. It is an interesting exercise. But if you tackle it, you will see that no matter
how you twist and turn trying to avoid the Cantor "hook", you are still caught on it. One
might say that any self-proclaimed "table of all reals" is hoist by its own petard.
        The repeatability of Cantor's diagonal method is similar to the repeatability of the
Tortoise's diabolic method for breaking the Crab's phonographs, one by one, as they got
more and more "hi-fi" and-at least so the Crab hoped-more "Perfect". This method
involves constructing, for each phonograph, a particular song which that phonograph
cannot reproduce. It is not a coincidence that Cantor's trick and the Tortoise's trick share
this curious repeatability; indeed, the Contracrostipunctus might well have been named
"Cantorcrostipunctus" instead. Moreover, as the Tortoise subtly hinted to the innocent
Achilles, the events in the Contracrostipunctus are a paraphrase of the construction which
Gödel used in proving his Incompleteness Theorem; it follows that the Gödel
construction is also very much like a diagonal construction. This will become quite
apparent in the next two Chapters.

                                From BlooP to FlooP
We have now defined the class of primitive recursive functions and primitive recursive
properties of natural numbers by means of programs written in the language BlooP. We
have also shown that BlooP doesn't capture all the functions of natural numbers which we
can define in words. We even constructed an "unBlooPable" function, Bluediag [N], by
Cantor's diagonal method. What is it about BlooP that makes Bluediag unrepresentable in
it? How could BlooP be improved so that Bluediag became representable?
        BlooP's defining feature was the boundedness of its loops. What if we drop that
requirement on loops, and invent a second language, called "FlooP" (`F' for "free")?
FlooP will be identical to BlooP except in one respect: we may have loops without
ceilings, as well as loops with ceilings (although the only reason one would include a
ceiling when writing a loop-statement in FlooP would be for the sake of elegance). These
new loops will be called MU-LOOPS. This follows the convention of mathematical
logic, in which "free" searches (searches without bounds) are usually indicated by a
symbol called a "µ-operator" (mu-operator). Thus, loop statements in FlooP may look
like this:

                             MU-LOOP:
                             BLOCK n: BEGIN
                                      .
                                      .
                             BLOCK n: END




BlooP and FlooP and GlooP                                                                431

   This feature will allow us to write tests in FlooP for such properties as wondrousness
and the Tortoise property-tests which we did not know how to program in BlooP because
of the potential open-endedness of the searches involved. I shall leave it to interested
readers to write a FlooP test for wondrousness which does the following things:

  (1) If its input, N, is wondrous, the program halts and gives the answer YES.
  (2) If N is unwondrous, but causes a closed cycle other than 1-4-2-1-4-2-1- ... , the
     program halts and gives the answer NO.
  (3) If N is unwondrous, and causes an "endlessly rising progression", the program
     never halts. This is FlooP's way of answering by not answering. FlooP's
     nonanswer bears a strange resemblance to Joshu's nonanswer "MU".

The irony of case 3 is that OUTPUT always has the value NO, but it is always
inaccessible, since the program is still grinding away. That troublesome third alternative
is the price that we must pay for the right to write free loops. In all FlooP programs
incorporating the MU-LOOP option, nontermination will always be one theoretical
alternative. Of course there will be many FlooP programs which actually terminate for all
possible input values. For instance, as I mentioned earlier, it is suspected by most people
who have studied wondrousness that a FlooP program such as suggested above will
always terminate, and moreover with the answer YES each time.

             Terminating and Nonterminating FlooP Programs
It would seem extremely desirable to be able to separate FlooP procedures into two
classes: terminators and nonterminators. A terminator will eventually halt no matter what
its input, despite the "MU-ness" of its loops. A nonterminator will go on and on forever,
for at least one choice of input. If we could always tell, by some kind of complicated
inspection of a FlooP program, to which class it belonged, there would be some
remarkable repercussions (as we shall shortly see). Needless to say, the operation of
class-checking would itself have to be a terminating operation-otherwise
one would gain nothing!

                                  Turing's Trickery
The idea springs to mind that we might let a BlooP procedure do the inspection. But
BlooP procedures only accept numerical input, not programs! However, we can get
around that ... by coding programs into numbers! This sly trick is just Gödel-numbering
in another of its many




BlooP and FlooP and GlooP                                                              432

a very long Gödel number. For instance, the shortest BlooP function (which is also a
terminating FlooP program)

            DEFINE PROCEDURE "A" [B]:
            BLOCK 0: BEGIN
            BLOCK 0: END.

-would get the Godel number partially shown below:

      904, 905, 906, 909, 914, 905           905, 914.904, 955,
      DEFINE                                        END.

      Now our scheme would be to write a BlooP test called TERMINATOR? which
says YES if its input number codes for a terminating FlooP program, NO if not. This way
we could hand the task over to a machine and with luck, distinguish terminators from
nonterminators. However, an ingenious argument given by Alan Turing shows that no
BlooP program can make this distinction infallibly. The trick is actually much the same
as Gödel’s trick, and therefore closely related to the Cantor diagonal trick. We shall not
give it here-suffice it to say that the idea is to feed the termination tester its own Godel
number. This is not so simple, however, for it is like trying to quote an entire sentence
inside itself. You have to quote the quote, and so forth; it seems to lead to an infinite
regress. However, Turing figured out a trick for feeding a program its own Godel
number. A solution to the same problem in a different context will be presented next
Chapter. In the present Chapter, we shall take a different route to the same goal, which is
namely to prove that a termination tester is impossible. For readers who wish to see an
elegant and simple presentation of the Turing approach, I recommend the article by
Hoare and Allison, mentioned in the Bibliography.

                   A Termination Tester Would Be Magical
Before we destroy the notion, let us delineate just why having a termination tester would
be a remarkable thing. In a sense, it would be like having a magical dowsing rod which
could solve all problems of number theory in one swell FlooP. Suppose, for instance, that
we wished to know if the Goldbach Variation is a true conjecture or not. That is, do all
numbers have the Tortoise property? We would begin by writing a FlooP test called
TORTOISE? which checks whether its input has the Tortoise property. Now the defect
of this procedure-namely that it doesn't terminate if the Tortoise property is absent-here
turns into a virtue! For now we run the termination tester on the procedure TORTOISE?.
If it says YES, that means that TORTOISE? terminates for all values of its input-in other
words, all numbers have the Tortoise property. If it says NO, then we know there exists a
number which has the Achilles property. The irony is that we never actually use the
program TORTOISE at all—we just inspect it.
         This idea of solving any problem in number theory by coding it into a




BlooP and FlooP and GlooP                                                               433

program and then waving a termination tester over the program is not unlike the idea of
testing a khan for genuineness by coding it into a folded string and then running a test for
Buddha-nature on the string instead. As
Achilles suggested, perhaps the desired information lies "closer to the surface" in one
representation than in another.

                 Pool F, Index Numbers, and Green Programs

Well, enough daydreaming. How can we prove that the termination tester is impossible?
Our argument for its impossibility will hinge on trying to apply the diagonal argument to
FlooP, just as we did to B1ooP. We shall see that there are some subtle and crucial
differences between the two cases.
        As we did for BlooP, imagine the pool of all FlooP programs. We shall call it
"Pool F". Then perform the same three filtering operations on Pool F, so that you get, in
the end:

   A complete pool of all call-less FlooP programs which calculate functions of
   exactly one input parameter.

Let us call these special FlooP-programs Green Programs (since they may go forever).
        Now just as we assigned index numbers to all Blue Programs, we can assign
index numbers to Green Programs, by ordering them in a catalogue, each volume of
which contains all Green Programs of a fixed length, arranged in alphabetical order.
So far, the carry-over from BlooP to FlooP has been straightforward. Now let us see if we
can also carry over the last part: the diagonal trick. What if we try to define a diagonal
function?

                      Greendiag [N] = 1 + Greenprogram{\#N} [N]

Suddenly, there is a snag: this function Greendiag [N] may not have a well-defined
output value for all input values N. This is simply because we have not filtered out the
nonterminator programs from Pool F, and therefore we have no guarantee that we can
calculate Greendiag [N] for all values of N. Sometimes we may enter calculations which
never terminate. And the diagonal argument cannot be carried through in such a case, for
it depends on the diagonal function having a value for all possible inputs.

               The Termination Tester Gives Us Red Programs

To remedy this, we would have to make use of a termination tester, if one existed. So let
us deliberately introduce the shaky assumption that one exists, and let us use it as our
fourth filter. We run down the list of Green Programs, eliminating one by one all
nonterminators, so that in the end we are left with:




BlooP and FlooP and GlooP                                                               434

   A complete pool of all call-less FlooP programs which calculate functions of
   exactly one input parameter, and which terminate for all values of their input..

Let us call these special FlooP programs Red Programs (since they all must stop). Now,
the diagonal argument will go through. We define

                         Reddiag [N] = 1 + Redprogram(\#N} [N]

and in an exact parallel to Bluediag, we are forced to conclude that Reddiag [N] is a well-
defined, calculable function of one variable which is not in the catalogue of Red
Programs, and is hence not even calculable in the powerful language FlooP. Perhaps it is
time to move on to GlooP?

                                        GlooP ...
Yes, but what is GlooP? If FlooP is BlooP unchained, then GlooP must be FlooP
unchained. But how can you take the chains off twice\% How do you make a language
whose power transcends that of FlooP? In Reddiag, we have found a function whose
values we humans know how to calculate-the method of doing so has been explicitly
described in English-but which seemingly cannot be programmed in the language FlooP.
This is a serious dilemma because no one has ever found any more powerful computer
language than FlooP.
         Careful investigation into the power of computer languages has been carried out.
We need not do it ourselves; let it just be reported that there is a vast class of computer
languages all of which can be proven to have exactly the same expressive power as FlooP
does, in this sense: any calculation which can be programmed in any one of the languages
can be programmed in them all. The curious thing is that almost any sensible attempt at
designing a computer language ends up by creating a member of this class-which is to
say, a language of power equal to that of FlooP. It takes some doing to invent a
reasonably interesting computer language which is weaker than those in this class. BlooP
is, of course, an example of a weaker language, but it is the exception rather than the rule.
The point is that there are some extremely natural ways to go about inventing algorithmic
languages; and different people, following independent routes, usually wind up creating
equivalent languages, with the only difference being style, rather than power.

                                      ... Is a Myth
In fact, it is widely believed that there cannot be any more powerful -language for
describing calculations than languages that are equivalent to FlooP. This hypothesis was
formulated in the 1930's by two people, independently of each other: Alan Turing—about
whom we shall say more later—and Alonzo Church, one of the eminent logicians of this
century. It




BlooP and FlooP and GlooP                                                                435

is called the Church-Turing Thesis. If we accept the CT-Thesis, we have to conclude that
"GlooP" is a myth-there are no restrictions to remove in FlooP, no ways to increase its
power by "unshackling" it, as we did BlooP.
         This puts us in the uncomfortable position of asserting that people can calculate
Reddiag [N] for any value of N, but there is no way to program a computer to do so. For,
if it could be done at all, it could be done in FlooP-and by construction, it can't be done in
FlooP. This conclusion is so peculiar that it should cause us to investigate very carefully
the pillars on which it rests. And one of them, you will recall, was our shaky assumption
that there is a decision procedure which can tell terminating from nonterminating FlooP
programs. The idea of such a decision procedure already seemed suspect, when we saw
that its existence would allow all problems of number theory to be solved in a uniform
way. Now we have double the reason for believing that any termination test is a myth-
that there is no way to put FlooP programs in a centrifuge and separate out the
terminators from the nonterminators.
         Skeptics might maintain that this is nothing like a rigorous proof that such a
termination test doesn't exist. That is a valid objection; however, the Turing approach
demonstrates more rigorously that no computer program can be written in a language of
the FlooP class which can perform a termination test on all FlooP programs.

                             The Church-Turing Thesis
Let us come back briefly to the Church-Turing Thesis. We will talk about it-and
variations on it-in considerable detail in Chapter XVII; for now it will suffice to state it in
a couple of versions, and postpone discussion of its merits and meanings until then. Here,
then, are three related ways to state the CT-Thesis:

   (1) What is human-computable is machine-computable.
   (2) What is machine-computable is FlooP-computable.
   (3) What is human-computable is FlooP-computable
         (i.e., general or partial recursive).

                  Terminology: General and Partial Recursive
We have made a rather broad survey, in this Chapter, of some notions from number
theory and their relations to the theory of computable functions. It is a very wide and
flourishing field, an intriguing blend of computer science and modern mathematics. We
should not conclude this Chapter without introducing the standard terminology for the
notions we have been dealing with.
        As has already been mentioned, “BlooP-computable” is synonymous with
“primitive recursive”. Now FlooP computable functions can be di-




BlooP and FlooP and GlooP                                                                  436

vided into two realms: (1) those which are computable by terminating FlooP programs:
these are said to be general recursive; and (2) those which are computable only by
nonterminating FlooP programs: these are said to be partial recursive. (Similarly for
predicates.) People often just say "recursive" when they mean "general recursive".

                                 The Power of TNT
It is interesting that TNT is so powerful that not only are all primitive recursive
predicates represented, but moreover all general recursive predicates are represented. We
shall not prove either of these facts, because such proofs would be superfluous to our
aim, which is to show that TNT is incomplete. If TNT could not represent some
primitive or general recursive predicates, then it would be incomplete in an uninteresting
way-so we might as well assume that it can, and then show that it is incomplete in an
interesting way.




BlooP and FlooP and GlooP                                                             437

